\documentclass[10pt,a4paper]{amsart}
\usepackage[margin=19mm]{geometry}
\usepackage{amsmath,amssymb,mathtools}
\usepackage[scr=rsfs,cal=euler]{mathalfa}
\usepackage{xfrac}
\usepackage[unicode,hidelinks,bookmarks=false]{hyperref}
\newcommand{\ie}{i.e.\ }
\newcommand{\cf}{cf.\ }
\newcommand{\resp}{respectively\ }
\newcommand{\Subsection}{\subsection}
\providecommand{\nobreakdash}{\nobreak\hbox{-}}
\setlength{\emergencystretch}{2em}
\multlinegap0pt
\usepackage{bm}
\usepackage{bbm}
\usepackage{dsfont}
\usepackage{xspace}
\usepackage{cancel}
\usepackage{mathtools}
\usepackage{amsthm}
\makeatletter
\@ifundefined{Theorem}{\newtheorem{Theorem}{Theorem}}{}
\@ifundefined{Proposition}{\newtheorem{Proposition}[Theorem]{Proposition}}{}
\makeatother
\newcommand\la{\lambda}
\title[On Polyhedral Formulas for Kirillov-Reshetikhin Modules]{On Polyhedral Formulas for Kirillov-Reshetikhin Modules}
\author{Chul-hee Lee}
\date{Source: arXiv:1901.00104v3 (CC BY 4.0)}
\begin{document}
\maketitle

\noindent\emph{Source and licence.} On Polyhedral Formulas for Kirillov-Reshetikhin Modules. Version arXiv:1901.00104v3 (CC BY 4.0), distributed under the Creative Commons Attribution 4.0 International licence, as recorded in the arXiv licence metadata for this exact version and shown on the abstract page https://arxiv.org/abs/1901.00104v3. Full rights notice: licenses/arxiv-1901.00104-CC-BY-4.0.txt.\par\smallskip

\noindent\textit{Editorial scope.} Counted unit: the Proposition whose source label is \texttt{prop:pfd} in the article (source0.tex lines 451--463, source label \texttt{prop:pfd}), delivered together with the article's own complete original proof of it (source0.tex lines 465--500, one \texttt{proof} environment of 36 lines). No mathematics is written, repaired or re-derived: statement and proof are the article's own text line for line. Required context retained uncounted: none. Every definition, display and label of those lines is retained unchanged. Omitted from this package and not redistributed: 1--450, 464--464, 501--919. Nothing inside a retained excerpt is omitted or altered: every retained body is delivered exactly as the article's lines are written, so no omission receipt of any kind accompanies this package. Citations: the retained excerpts carry no citation command at all, so no bibliography entry of the article is transcribed and no reference is redistributed. Reference adaptation: none was needed and none was made; every reference inside the delivered unit resolves inside it, so no unresolved reference or citation is produced. Printed slips kept literal: no printed word, symbol, hypothesis or number of the counted statement or of its proof was altered. Layout and class: the article's own class is \texttt{sigma} with packages including ifpdf, eucal, enumitem; the wrapper is the lane's pinned \texttt{amsart} editorial wrapper at 10pt on A4, which restates the theorem-like environments the retained text needs and adds the article's own macro definitions that the retained text needs (\texttt{\textbackslash la}), with extra wrapper packages (bm, bbm, dsfont, xspace, cancel, mathtools). The delivered numbering is the unit's own. Assets: no retained span contains a figure environment, a graphics-inclusion command, a PDF-page inclusion, a table or a TikZ picture; no image file is copied into this package, the package declares no asset dependency and no image file is redistributed with it. Licence: version arXiv:1901.00104v3 of this article is distributed under the Creative Commons Attribution 4.0 International licence, as recorded in the arXiv licence metadata for this exact version and shown on the abstract page https://arxiv.org/abs/1901.00104v3. A full rights notice is delivered at licenses/arxiv-1901.00104-CC-BY-4.0.txt. Characters: every retained excerpt is pure ASCII, so the delivered engine, XeLaTeX with its default Latin Modern text font, reports no missing character.
\begin{Proposition}\label{prop:pfd}
We have
\[
\mathcal{P}^{(2)}(t) =
\frac{1}{\Delta}\sum_{w\in W}(-1)^{\ell(w)}{\rm e}^{w(\rho)}
\left(
\frac{w\bigl(d^{(2)}_{0}\bigr)}{1-t}
+\frac{w\bigl(d^{(2)}_{\la_1}\bigr)}{1-{\rm e}^{w(\la _1)}t}
+\frac{w\bigl(d^{(2)}_{\la_3}\bigr)}{1-{\rm e}^{w(\la _3)}t}
+\frac{w\bigl(d^{(2)}_{\la_4}\bigr)}{1-{\rm e}^{w(\la _4)}t}
\right).
\]
\end{Proposition}

\begin{proof}
It is sufficient to show that for $\la\in \{\la_2,\la_3,\la_4\}$,
\begin{equation}\label{sumEvanishing}
\sum_{w\in W_{\la}}(-1)^{\ell(w)}{\rm e}^{w(\rho)}w\bigl(e^{(2)}_{\la}\bigr)=0.
\end{equation}
We may use computers to verify this directly. Below, we will explain how to reduce the amount of calculation to check~\eqref{sumEvanishing}. While this reduction is not essential as long as we focus on type~$F_4$ whose Weyl group is manageable in size, it might be useful for treating a~similar vanishing sum over bigger groups in other types.

For $\la=\la_2=2\omega_3$, $W_{\la}=W_{\{1,2,4\}}$. The parabolic subgroup $W_{\{1,2\}}$ of $W_{\{1,2,4\}}$ satisfies
\[
\sum_{w\in W_{\{1,2\}}}(-1)^{\ell(w)}{\rm e}^{w(\rho)}w\bigl(e^{(2)}_{\la_2}\bigr)=0,
\]
which implies the vanishing of the sum over $W_{\{1,2,4\}}$ since
\eqref{sumEvanishing} can be written as
\[
\sum_{w'\in W'}(-1)^{\ell(w')}w'\left(\sum_{w\in W_{\{1,2\}}}(-1)^{\ell(w)}{\rm e}^{w(\rho)}w\bigl(e^{(2)}_{\la_2}\bigr)\right),
\]
where $W'$ is the set of minimal coset representatives for cosets in $W_{\{1,2,4\}}/W_{\{1,2\}}$.

Similarly, when $\la=\la_3=\omega_2$, we have $W_{\la}=W_{\{1,3,4\}}$ and the following sum over the parabolic subgroup $W_{\{1,3\}}$ vanishes:
\[
\sum_{w\in W_{\{1,3\}}}(-1)^{\ell(w)}{\rm e}^{w(\rho)}w\bigl(e^{(2)}_{\la_3}\bigr)=0.
\]

When $\la=\la_4=\omega_1$, we have not found any proper parabolic subgroup of $W_{\la}=W_{\{2,3,4\}}$, over which the sum vanishes. However, if we let
\[
E^{(2)}_{\la_4}:=\sum_{w\in W_{\{2,4\}}}(-1)^{\ell(w)}{\rm e}^{w(\rho)}w\bigl(e^{(2)}_{\la_4}\bigr),
\]
then the left-hand side of~\eqref{sumEvanishing} becomes
\[
\sum_{w\in W''}(-1)^{\ell(w)} w\bigl(E^{(2)}_{\la_4}\bigr),
\]
where $W''$ is the set of minimal coset representatives for cosets in $W_{\{2,3,4\}}/W_{\{2,4\}}$. The size of~$W''$ is 12, and it is possible to partition this set into 6 pairs of distinct elements so that the contribution from each pair to the above sum is zero. In other words, for each $w_1\in W''$, there exists $w_2\in W''$, $w_1\neq w_2$ such that
\[
(-1)^{\ell(w_1)} w_1\bigl(E^{(2)}_{\la_4}\bigr)+(-1)^{\ell(w_2)}w_2\bigl(E^{(2)}_{\la_4}\bigr)=0.\tag*{\qed}
\] \renewcommand{\qed}{}
\end{proof}

\end{document}
